\chapter{Основные объекты памяти}
\label{ch:definitions}

\section{Зачем вообще нужна отдельная глава понятий}

В технических текстах о памяти часто возникает неприятная ситуация: читатель
уже встретил ``уверенность'', ``активацию'', ``провенанс'', ``консолидацию'' и
``walker'', но ещё не получил компактной картины того, как эти слова
соотносятся друг с другом. Формально каждое из них потом где-то определяется,
однако понимание успевает распасться на локальные фрагменты.

Задача этой главы обратная: собрать в одном месте основные объекты, не
пытаясь здесь же полностью исчерпать все механизмы. Поэтому формулировки
ниже намеренно коротки. Их роль --- дать рабочий словарь и развести сущности,
которые особенно легко спутать.

\section{Объекты субстрата}

\begin{definition}[Граф памяти]
Граф памяти $G = (V, E)$ --- направленный динамический граф, в котором узлы
$V$ несут содержательные единицы памяти, а рёбра $E$ несут отношения между
ними. Граф называется типизированным, потому что и узлы, и рёбра имеют виды;
темпоральным, потому что память различает временные слои и возраст
содержания; динамическим, потому что операторы могут менять и узлы, и рёбра.
\end{definition}

\begin{definition}[Узел памяти]
Узел памяти $v \in V$ --- это единица содержания, с которой можно работать как
с объектом: наблюдение, факт, эпизод, гипотеза, предпочтение, паттерн
реакции, модель отношения, элемент самопредставления и так далее. Узел хранит
не только текст или embedding, но и поля состояния: тип узла, уверенность,
провенанс, временной слой, активацию, валентность, важность и служебные
флаги.
\end{definition}

\begin{intuition}[Почему узел --- это больше, чем ``запись'']
Если узел понимать только как текстовый фрагмент, память превращается в
коллекцию документов. Узел в этой архитектуре ближе к ``объекту мышления'':
он не просто что-то содержит, но ещё может быть активирован, усилен, ослаблен,
проверен, связан с посылками и включён в дальнейший вывод.
\end{intuition}

\begin{definition}[Тип узла]
\texttt{node\_kind} --- вид узла, определяющий, как именно его следует читать и
обрабатывать. В частности, именно тип узла задаёт смысл поля confidence. Для
наблюдаемого факта это степень доверия к истинности; для предпочтения ---
устойчивость оценки; для реактивного паттерна --- надёжность модели ответа;
для mood-summary это может быть вообще не числовая истина, а служебная
свёртка состояния.
\end{definition}

\begin{pitfall}[Почему тип узла не декоративен]
В памяти проекта сосуществуют по меньшей мере observation, claim, hypothesis,
episode, preference, reaction\_pattern, relationship\_state, boundary,
skill/procedure, self\_model и mood-summary. Эти сущности нельзя честно
измерять одной и той же шкалой: ``уверенность 0.8'', относящаяся к факту, и
``уверенность 0.8'', относящаяся к предпочтению, означают разные инженерные
вещи.
\end{pitfall}

\begin{definition}[Ребро]
Ребро $(u \to v)$ --- направленная связь между двумя узлами, несущая тип
отношения. Примеры: таксономическое включение, причинная поддержка,
противоречие, ассоциация, участие в одном эпизоде, отношение к одному
участнику взаимодействия.
\end{definition}

\begin{definition}[Два канала ребра]
На ребре живут два разных числовых канала:
\begin{itemize}
\item вес перехода $w_{uv} \ge 0$, который нужен обходу и говорит, насколько охотно маршрут пойдёт по этой связи;
\item знак $s_{uv} \in \{+1, -1\}$, который нужен эпистемике и говорит, усиливает ли эта связь убеждение или ослабляет его.
\end{itemize}
Эти каналы не следует смешивать: маршрут может часто ходить по ребру, которое
при этом вносит отрицательный вклад в уверенность, и наоборот.
\end{definition}

\begin{definition}[Провенанс (обоснование)]
Провенанс --- это информация о происхождении узла: из какого наблюдения,
события, источника или набора посылок он возник. Для выводных узлов
провенанс играет роль обоснования, то есть justification set. Важно, что
provenance --- это не полный трек обхода: маршрут walker'а шумен, может
храниться отдельно и сам по себе не является обоснованием. Именно provenance
делает возможной ретракцию: если опровергнута посылка, можно найти и её
следствия.
\end{definition}

\begin{definition}[Время: системное и содержательное]
Время в памяти удобно разводить на два канала. Системное время --- это
минимальный набор меток, нужный самим операторам: \texttt{layer} (сцена,
сессия, день, долгий паттерн) и, как правило, одна служебная метка вроде
времени кодирования или последнего доступа. Содержательное время --- это ответ
на вопрос ``когда это произошло'' внутри самого мира памяти. Оно может
моделироваться отдельными узлами и быть реконструктивным, а значит, и
ошибочным. Интервалы валидности (\texttt{start}/\texttt{end}) имеет смысл
вешать лишь на те факты, для которых это содержательно нужно, а не
механически на каждый узел.
\end{definition}

\section{Поля состояния}

\begin{definition}[Уверенность]
Уверенность $\conf(v) \in [-1, 1]$ --- знаковая степень убеждённости в
содержании узла. Положительный знак означает ``скорее да'', отрицательный ---
``скорее нет'', нуль соответствует незнанию или эпистемической
неопределённости. Такая запись удобна тем, что различает два принципиально
разных состояния: ``неверие'' и ``отсутствие знания''.
\end{definition}

\begin{equation}
\conf = 2P - 1,
\end{equation}

где $P$ --- эквивалентная вероятность истинности. Это соответствие полезно, но
в книге предпочтителен именно знаковый скаляр: он лучше согласуется с
противоречием и ретракцией. Практически удобно считать, что зона незнания
задаётся условием $|\conf(v)| < h$ с грубым ориентиром $h \approx 0{,}3$.
Именно такие узлы являются естественными кандидатами на
\texttt{need\_to\_discover}, пересмотр или удаление. Уверенность меняется
медленно: через пропагацию и ревизию убеждений, а не при каждом локальном
подсвечивании. Принципиальное ограничение: выводный узел не получает
$|\conf| = 1$.

\begin{pitfall}[Чего нельзя требовать от confidence]
Одно число не обязано кодировать всю эпистемическую жизнь узла. Оно не
заменяет тип узла, не заменяет провенанс и не заменяет внешние механизмы
неопределённости, такие как Beta-модель реакций или флаги \texttt{need\_to\_*}.
Поэтому вопрос ``почему бы не хранить всё в одном идеальном скаляре'' почти
всегда поставлен неверно.
\end{pitfall}

\begin{definition}[Активация]
Активация $a_v(t) \ge 0$ --- краткоживущая мера того, насколько узел
``подсвечен сейчас''. Активация относится к локальной динамике памяти, а не к
долговременной вере. Она может возрастать из-за запроса, обхода, совместной
активности соседей или внутренней заинтересованности системы.
\end{definition}

\begin{equation}
a(t + 1) = \lambda \tilde{P}^{\top} a(t) + s(t), \qquad \lambda \in (0, 1),
\end{equation}

где $\tilde{P}$ --- матрица переходов, а $s(t)$ --- внешний или внутренний
вброс активации. Эквивалентно об этом можно думать как об ожидаемой частоте
посещений узла роем в коротком окне. Активацию важно не путать ни с
уверенностью, ни со стационарным распределением $\pi$: первое отвечает за
убеждение, второе --- за долгосрочную геометрию потока.

\begin{definition}[Валентность и настроение]
Валентность --- аффективная окраска узла. Она отвечает не за то, насколько узел
истинный, а за то, как он переживается и как он взаимодействует с настроением.
В книге обычно предполагается двумерное описание ``валентность--возбуждение'',
но точная параметризация вторична по отношению к самой идее аффективной
несущей. Настроение \texttt{mood} --- отдельный низкоразмерный вектор
состояния, сцепленный с графом двусторонне, но не живущий внутри каждого узла
по отдельности.
\end{definition}

\begin{definition}[Значимость]
Значимость (\texttt{importance}, \texttt{salience}) --- мера того, насколько
объект памяти важен для дальнейшей судьбы системы. Значимость влияет на
консолидацию, забывание и приоритеты переработки. Она может частично зависеть
от эмоционального сдвига, от роли объекта в долгих паттернах и от структурной
центральности. В качестве дешёвого прокси значимости удобно использовать
величину изменения настроения $|\Delta \mood|$; более глубокий, но и более
дорогой сигнал --- ошибка предсказания.
\end{definition}

\section{Операторы над памятью}

\begin{definition}[Random Walker]
Random Walker --- блуждающий агент, который движется по графу по
вероятностному правилу, зависящему от структуры рёбер, релевантности,
текущего настроения, неопределённости и других факторов. В типовой форме:
\[
P(A \to B) \propto w_{AB} \cdot rel(B) \cdot mc(B) \cdot unc(B)^{\beta_u},
\]
где $w_{AB}$ --- вес перехода, $rel(B)$ --- релевантность узла текущему
запросу или сцене, $mc(B)$ --- mood-congruent bias, а $unc(B)^{\beta_u}$ ---
смещение к неопределённому содержанию в тех режимах, где это нужно. Один
walker даёт локальную траекторию; множество walker'ов даёт статистическую
картину потока по памяти.
\end{definition}

Функционал одного walker'а в минимуме таков: (i) перейти по вероятностному
правилу; (ii) с вероятностью $\alpha$ телепортироваться, то есть
перезапуститься из другого узла; (iii) повысить активацию посещённого узла;
(iv) по условию запустить операцию вроде retrieval, fusion или инспекции
\texttt{need\_to\_*}; (v) записать след как последовательность посещённых узлов.

\begin{definition}[Рой обходчиков]
Рой (\texttt{swarm}) --- набор walker'ов, работающих параллельно под общей
политикой. Именно рой, а не одиночный маршрут, лучше соответствует идее
распределённой активации памяти. На достаточной статистике следов рой даёт и
эмпирическую картину маршрутов по графу, что особенно важно для
экспериментов.
\end{definition}

\begin{definition}[Три политики обхода]
Один и тот же движок обхода должен уметь работать как минимум в трёх режимах:
retrieval под запрос, внутренняя жизнь / любопытство с повышенным интересом к
неопределённому и сон, где включается более широкое стохастическое покрытие с
телепортом. Эти режимы различаются не типом сущности, а параметрами той же
машины.
\end{definition}

\begin{definition}[Retrieval]
Retrieval --- извлечение содержательно релевантной части памяти. Здесь это не
один шаг, а двухходовая процедура: сначала находятся якоря быстрым индексом,
затем от якорей запускается ограниченное расширение окрестности посредством
обхода. Этот порядок принципиален: walker хорош для graph expansion, но плох
как механизм поиска входа. В результате retrieval зависит не только от
семантической близости запроса, но и от настроения, то есть допускает
mood-congruent recall.
\end{definition}

\begin{definition}[Fusion]
Fusion --- механизм рождения нового узла из согласованной конфигурации
совместно активированных посылок. Обычно результат fusion --- гипотеза или
обобщение с ограниченной уверенностью и явным провенансом. Для выводных
узлов принципиально: результат fusion не стартует как ``окончательный факт'',
а сохраняет провизорный статус до дополнительной проверки.
\end{definition}

\begin{definition}[Консолидация]
Консолидация --- процесс перехода от множества частных следов к более
устойчивым структурам памяти. Она может включать объединение дублей,
summary-overlay, рождение обобщений и более позднее огрубление графа. Важно,
что консолидация --- не один жёсткий оператор, а семейство подопераций:
fusion/вывод, summary-overlay, dedup/merge и только на более позднем этапе
разрушающее огрубление.
\end{definition}

\begin{definition}[Забывание]
Забывание --- ослабление или удаление тех элементов памяти, которые не
поддерживаются ни недавним использованием, ни значимостью, ни более глубокой
структурной ролью. В этой архитектуре забывание не противоположно
консолидации, а образует с ней единую динамику перераспределения памяти во
времени. На базовом уровне его удобно связывать с частотой и свежестью
доступа, то есть с медленно меняющейся базовой активацией.
\end{definition}

\begin{definition}[Сон]
Сон --- фоновый режим памяти, в котором обход становится более стохастическим
и более широким, чем в дневном retrieval. В этом режиме происходят скраб,
пересмотр противоречий, консолидация, случайные инспекции и иногда
неожиданные дальние fusion-связки. Тот же режим удобно использовать и для
подтягивания внешнего материала по флагам \texttt{need\_to\_discover}.
\end{definition}

\section{Сцепленные с памятью процессы}

\begin{definition}[Настроение]
Настроение \texttt{mood} --- отдельный вектор состояния, который не хранится
внутри каждого узла, но двусторонне сцеплен с графом. Настроение влияет на
маршруты активации и на форму выражения; граф, в свою очередь, поставляет
настроению материал и обратную связь.
\end{definition}

\begin{definition}[Личность]
Личность --- не таблица жёстко прошитых весов, а способ, которым мотивация,
настроение, эпистемическая осторожность и поведенческие приоритеты смещают
работу памяти и её последующее выражение. Внутри архитектуры личность
проявляется не как отдельный \texttt{store}, а как \texttt{bias} во время обхода
и как стиль на выходе.
\end{definition}

\begin{definition}[Голос]
Голос --- языковой слой, который превращает внутренний контекст памяти в
конкретную реплику. Голос не хранит сами факты вместо памяти, а композирует
высказывание из того, что память уже подняла и подготовила. Часть характера
живёт именно здесь; факты, \texttt{provenance} и эпистемические отношения
живут в графе.
\end{definition}

\section{Стартовое состояние и инженерная перспектива}

\begin{definition}[Эмбриональный граф]
Эмбриональный граф --- стартовое заполнение памяти из внешних баз знаний и
здравого смысла, прежде всего в духе ConceptNet + ATOMIC. Это не готовая
личность и не законченная память, а начальная, нейтральная по характеру среда,
в которой могут запускаться собственные механизмы retrieval, fusion,
консолидации и ревизии.
\end{definition}

\begin{definition}[Stage 1]
Stage 1 --- первый рабочий вертикальный срез системы, в котором теоретическая
модель памяти доводится до минимально наблюдаемой инженерной машины: с
\texttt{event log}, \texttt{write-path}, \texttt{typed node store}, retrieval,
\texttt{prompt packing}, инспектором памяти и синтетическими тестами
(глава 20).
\end{definition}

\begin{pitfall}[Почему не вводится более богатая эпистемическая тройка]
В исходной архитектуре сознательно удерживается знаковый скаляр
$\conf \in [-1, 1]$, а не полный кортеж belief/disbelief/uncertainty. Это
решение не объявляется окончательно истинным: более богатая субъективная
логика остаётся допустимой альтернативой, если разнесённые сигналы
неопределённости окажутся неудобными. Но на текущем этапе roadmap делает
ставку именно на скаляр confidence, provenance, node\_kind и дополнительные
специальные поля неопределённости там, где они реально нужны.
\end{pitfall}

\section{Что должно остаться после этой главы}

После этой главы полезно удерживать три различения.

\begin{enumerate}
\item Узел не равен строке текста. Это содержательный объект с историей и состоянием.
\item Активация не равна уверенности. Первая отвечает за фокус сейчас, вторая --- за долговременное убеждение.
\item Оператор не равен полю хранения. Retrieval, fusion, консолидация и сон --- это процессы над графом, а не разные виды записей в одном и том же столбце.
\end{enumerate}

На этом словарь готов. Следующая глава уже может говорить не об именах
объектов, а о самом субстрате памяти как инженерной конструкции: какие поля
действительно нужны графу и почему именно они.
